\documentclass[10pt,a4paper]{amsart}
\usepackage[margin=19mm]{geometry}
\usepackage{amsmath,amssymb,mathtools}
\usepackage[scr=rsfs,cal=euler]{mathalfa}
\usepackage{xfrac}
\usepackage[unicode,hidelinks,bookmarks=false]{hyperref}
\newcommand{\ie}{i.e.\ }
\newcommand{\cf}{cf.\ }
\newcommand{\resp}{respectively\ }
\newcommand{\Subsection}{\subsection}
\providecommand{\nobreakdash}{\nobreak\hbox{-}}
\setlength{\emergencystretch}{2em}
\multlinegap0pt
\usepackage{bm}
\usepackage{bbm}
\usepackage{dsfont}
\usepackage{xspace}
\usepackage{cancel}
\usepackage{mathtools}
\usepackage[shortlabels]{enumitem}
\usepackage{amsthm}
\makeatletter
\@ifundefined{theorem}{\newtheorem{theorem}{Theorem}}{}
\@ifundefined{definition}{\newtheorem{definition}[theorem]{Definition}}{}
\makeatother
\newcommand\ord{\mbox{ord}}
\renewcommand\ge\geqslant
\renewcommand\geq\geqslant
\renewcommand\a{\alpha}
\renewcommand\le\leqslant
\renewcommand\leq\leqslant
\renewcommand\ge\geqslant
\renewcommand\geq\geqslant
\renewcommand\a{\alpha}
\title[Reconfiguration of Nowhere-zero Flows]{Reconfiguration of Nowhere-zero Flows}
\author{Daniel W. Cranston, Jiaao Li, Bo Su, Zhouningxin Wang,, Ningyan Xu}
\date{Source: arXiv:2606.24685v1 (CC BY 4.0)}
\begin{document}
\maketitle

\noindent\emph{Source and licence.} Reconfiguration of Nowhere-zero Flows. Version arXiv:2606.24685v1 (CC BY 4.0), distributed under the Creative Commons Attribution 4.0 International licence, as recorded in the arXiv licence metadata for this exact version and shown on the abstract page https://arxiv.org/abs/2606.24685v1. Full rights notice: licenses/arxiv-2606.24685-CC-BY-4.0.txt.\par\smallskip

\noindent\textit{Editorial scope.} Counted unit: the Theorem whose source label is \texttt{thm:universal-FRC} in the article (source0.tex lines 1139--1145, source label \texttt{thm:universal-FRC}), delivered together with the article's own complete original proof of it (source0.tex lines 1147--1198, one \texttt{proof} environment of 52 lines). No mathematics is written, repaired or re-derived: statement and proof are the article's own text line for line. Required context retained uncounted: $A$-FRC-def (lines 1033--1045); FRC-1 (lines 1037--1044); FRC-2 (lines 1037--1044); FRC-3 (lines 1037--1044). Every definition, display and label of those lines is retained unchanged. Omitted from this package and not redistributed: 1--1032, 1045--1138, 1146--1146, 1199--1788. Nothing inside a retained excerpt is omitted or altered: every retained body is delivered exactly as the article's lines are written, so no omission receipt of any kind accompanies this package. Citations: the retained excerpts carry no citation command at all, so no bibliography entry of the article is transcribed and no reference is redistributed. Reference adaptation: none was needed and none was made; every reference inside the delivered unit resolves inside it, so no unresolved reference or citation is produced. Printed slips kept literal: no printed word, symbol, hypothesis or number of the counted statement or of its proof was altered. Layout and class: the article's own class is \texttt{article} with packages including babel, geometry, indentfirst, hyperref, bm, tikz, caption, amsmath, amssymb, amsthm, mathtools, braket, enumitem, booktabs; the wrapper is the lane's pinned \texttt{amsart} editorial wrapper at 10pt on A4, which restates the theorem-like environments the retained text needs and adds the article's own macro definitions that the retained text needs (\texttt{\textbackslash ge}, \texttt{\textbackslash le}, \texttt{\textbackslash ord}), with extra wrapper packages (bm, bbm, dsfont, xspace, cancel, mathtools, enumitem). The delivered numbering is the unit's own. Assets: no retained span contains a figure environment, a graphics-inclusion command, a PDF-page inclusion, a table or a TikZ picture; no image file is copied into this package, the package declares no asset dependency and no image file is redistributed with it. Licence: version arXiv:2606.24685v1 of this article is distributed under the Creative Commons Attribution 4.0 International licence, as recorded in the arXiv licence metadata for this exact version and shown on the abstract page https://arxiv.org/abs/2606.24685v1. A full rights notice is delivered at licenses/arxiv-2606.24685-CC-BY-4.0.txt. Characters: every retained excerpt is pure ASCII, so the delivered engine, XeLaTeX with its default Latin Modern text font, reports no missing character.
\begin{definition}
\label{$A$-FRC-def}
Let $H$ be a graph and $A$ be an abelian group.
We say that $H$ is \emph{$A$-flow-reconfiguration-contractible}, or simply \emph{$A$-FRC}, if the following 3 conditions hold:
\begin{enumerate}[label=(\roman*)]
\item\label{FRC-1}
The graph $H$ is $A$-connected.
\item\label{FRC-2}
For every $A$-boundary $\beta$ of $H$, the graph $H$ is $\beta$-flow-connected.
\item\label{FRC-3}
For any two $A$-boundaries $\beta_1$ and $\beta_2$ that differ only at two vertices $v_1$ and $v_2$, the graph $H$ is $(\beta_1,\beta_2)$-flow-convertible.
\end{enumerate}
\end{definition}

\begin{theorem}\label{thm:universal-FRC}
Let $A$ be a finite abelian group, 
and fix an integer $n \ge 2$. 
For each element $\gamma \in A \setminus \{0_A\}$, denote its order by $\ord(\gamma)$. The \emph{structural cost} $M(\gamma)$ is given by 
\[ M(\gamma):= \frac{|A|}{\ord(\gamma)} \left\lceil \frac{2 \cdot \ord(\gamma)}{3} \right\rceil. \]
The cycle $C_n$ is $A$-FRC if and only if $M(\gamma) \ge n + 1$ for every $\gamma \in A \setminus \{0_A\}$.
\end{theorem}

\begin{proof}
We begin by algebraically
parameterizing all flows on a cycle $C_n$. We identify all orientations of $C_n$ with a fixed cyclic orientation $D$ by replacing the value on an oppositely oriented edge by its inverse in $A$. For each $A$-boundary $\beta$, a $\beta$-flow $(D, f)$ on $C_n$ is uniquely determined by the flow value $a = f(e_1)$ on an arbitrary reference edge $e_1$. By flow conservation, the flow on each other edge $e \in E(C_n)$ can be written as $f(e) = a + c_e$, where $c_e \in A$ is a constant offset determined entirely by $\beta$.
The nowhere-zero condition requires that $a+c_e \neq 0_A$ for every edge $e\in E(C_n)$; equivalently, the initial value $a$ must avoid the forbidden values $-c_e$ for all $e\in E(C_n)$.

We first prove the sufficiency. Assume $M(\gamma) \ge n + 1$ for all $\gamma \neq 0_A$. 
Let $p:=\ord(\gamma)$.
Now $M(\gamma) = \frac{|A|}{p}\lceil \frac{2p}{3} \rceil \le \frac{|A|}{p} \cdot p = |A|$. Thus, our hypothesis gives $|A| \ge n+1$.

We verify the three conditions of Definition~\ref{$A$-FRC-def}. For each $A$-boundary $\beta$, the nowhere-zero condition forbids at most $n$ initial values (one for each edge). Since $|A| \ge n+1 > n$, there is at least one valid nowhere-zero $\beta$-flow; hence, $C_n$ is $A$-connected (so satisfies Definition~\ref{$A$-FRC-def}\ref{FRC-1}). Furthermore, every two $\beta$-flows differ by a constant flow along the entire cycle. Thus, their difference is a valid cycle flow, meaning the reconfiguration graph $\mathcal{F}(C_n, \beta)$ is connected, fulfilling Definition~\ref{$A$-FRC-def}\ref{FRC-2}.

To verify that $C_n$ is
$(\beta_1, \beta_2)$-convertible (Definition~\ref{$A$-FRC-def}\ref{FRC-3}), let $\beta_1$ and $\beta_2$ differ by $\gamma \in A \setminus \{0_A\}$ at vertices $u$ and $w$. Let $P_1$ and $P_2$ be the two paths on $C_n$ connecting $u$ and $w$. Note that $|E(P_1)| + |E(P_2)| = n$. Let $X$ be the set of forbidden initial flow values associated with $P_1$ under $\beta_1$ (i.e., $X = \{-c_e \mid e \in E(P_1)\}$), and $Y$ be the forbidden set for $P_2$. Clearly, $|X| \le |E(P_1)|$ and $|Y| \le |E(P_2)|$, yielding $|X| + |Y| \le n$. 

Let $Z:=A\setminus (X\cup Y)$.
To convert $\beta_1$ to $\beta_2$, we must either add $\gamma$ to $P_1$ or subtract $\gamma$ from $P_2$. Suppose  that instead both options fail to produce a nowhere-zero $\beta_2$-flow for every valid initial value $z \in Z$. If adding $\gamma$ to $P_1$ fails, then the shifted values must hit the forbidden set of $P_1$, implying $Z + \gamma \subseteq X$. If subtracting $\gamma$ from $P_2$ fails, then $Z - \gamma \subseteq Y$. 

Note that the element $\gamma$ generates a cyclic subgroup $\langle \gamma \rangle$ of order $p$,
which partitions $A$ into $|A|/p$ disjoint cosets. For any such coset $A_i$, let $Z_i:= Z \cap A_i$, $X_i:= X \cap A_i$, and $Y_i:= Y \cap A_i$. The failure conditions apply locally: $Z_i + \gamma \subseteq X_i$ and $Z_i - \gamma \subseteq Y_i$. Let $z_i:= |Z_i|$, $x_i:= |X_i|$, and $y_i:= |Y_i|$. These subset inclusions force $x_i \ge z_i$ and $y_i \ge z_i$, yielding $x_i + y_i \ge 2z_i$. Moreover, since $Z_i$ is the complement of $X_i \cup Y_i$ within the coset, we have $|X_i \cup Y_i| = p - z_i$, forcing $x_i + y_i \ge p - z_i$. Thus, for every coset, we have 
\[ x_i + y_i \ge \max(2z_i, p - z_i). \]
The lower bound for this maximum over all integers $z_i \ge 0$ occurs near the real-valued intersection $2z_i = p - z_i$, meaning $z_i \approx p/3$. Letting $z_i:= \lfloor p/3 \rfloor$ yields the exact minimum $\min(x_i + y_i) = \lceil 2p/3 \rceil$. Summing this minimum bound over all $|A|/p$ cosets gives:
\[ |X| + |Y| \ge \sum_{i=1}^{|A|/p} \left\lceil \frac{2p}{3} \right\rceil = \frac{|A|}{p} \left\lceil \frac{2p}{3} \right\rceil = M(\gamma). \]
However, this implies $M(\gamma) \le |X| + |Y| \le n$, which directly contradicts our initial hypothesis that $M(\gamma) \ge n+1$. Therefore, at least one conversion shift must succeed for some valid nowhere-zero flow, proving that $C_n$ is $A$-FRC.
\medskip

We now prove the necessity by an explicit construction. Suppose that there exists an element $\gamma \in A \setminus \{0_A\}$ such that $M(\gamma) \le n$. We will construct a specific $A$-boundary $\beta_1$, and assign a target boundary $\beta_2$ that differs from $\beta_1$ by exactly $\gamma$ at two distinct vertices. We will demonstrate that either no valid $\beta_1$-flow exists, or that no valid $\beta_1$-flow can be successfully reconfigured into any valid $\beta_2$-flow. In either case, this violates the conditions of Definition~\ref{$A$-FRC-def}, proving that $C_n$ is not $A$-FRC. 

Let $p := \ord(\gamma)$. In each coset $A_i$ of the cyclic subgroup $\langle \gamma \rangle$, we can select sets $X_i, Y_i \subseteq A_i$ satisfying $|X_i| + |Y_i| = \lceil 2p/3 \rceil$ such that for any element $z \in A_i \setminus (X_i \cup Y_i)$, we have $z + \gamma \in X_i$ and $z - \gamma \in Y_i$. 

To see this combinatorial fact explicitly, we identify $A_i$ with the cyclic sequence $(v_0, v_1, \dots, v_{p-1})$, where $v_{j+1} = v_j + \gamma$ (indices taken modulo $p$). We can systematically place valid initial values by defining the set $Z_i = \{v_{3m} \mid 0 \le m < \lfloor p/3 \rfloor\}$, which has size $\lfloor p/3 \rfloor$. The required conditions dictate that the right neighbors $v_{3m+1}$ must belong to $X_i$ and the left neighbors $v_{3m-1}$ must belong to $Y_i$. Since these neighbor sets are mutually disjoint for $m < \lfloor p/3 \rfloor$, we strictly assign them to $X_i$ and $Y_i$, respectively. Any remaining unassigned elements in $A_i$ can be arbitrarily partitioned into $X_i$ and $Y_i$. This construction leaves exactly $\lfloor p/3 \rfloor$ elements in $Z_i = A_i \setminus (X_i \cup Y_i)$, ensuring that $|X_i| + |Y_i| = p - \lfloor p/3 \rfloor = \lceil 2p/3 \rceil$.

Taking the union across all $|A|/p$ cosets, we obtain global forbidden sets $X = \bigcup X_i$ and $Y = \bigcup Y_i$, with a total combined size of $|X| + |Y| = M(\gamma)$.
%
Because $M(\gamma) \le n$, the cycle $C_n$ has enough edges to accommodate these constraints. We partition $C_n$ into two paths, $P_1$ and $P_2$, both with endpoints $u$ and $w$, such that $|E(P_1)| \ge |X|$ and $|E(P_2)| \ge |Y|$. We then directly assign the offset $c_e$ to each edge $e$: for edges in $P_1$, we injectively assign the value $-x$ for every $x \in X$, and assign arbitrary values from $-X$ to any remaining edges. We apply the identical procedure for $P_2$, assigning $-y$ for every $y \in Y$. 

With all $c_e$ fixed, we define the boundary $\beta_1(v) = c_{e_{out}} - c_{e_{in}}$ for every vertex $v$, where $e_{in}$ and $e_{out}$ are the incoming and outgoing edges incident to $v$ under the chosen orientation $D$. This definition ensures that $\sum_{v \in V(C_n)} \beta_1(v) = 0_A$, making $\beta_1$ a valid boundary. 
Under this specific boundary $\beta_1$, if we choose an initial flow value $a = f_1(e_1)$, the flow on any arbitrary edge $e \in E(C_n)$ is exactly $f_1(e) = a + c_e$. To ensure $f_1$ is a valid nowhere-zero flow, the initial value $a$ must avoid $-c_e$ for all edges. By our assignment, the set of all such forbidden values is exactly $X \cup Y$. Thus, the set of valid initial values is precisely $Z = A \setminus (X \cup Y)$.

At this point, we define the target boundary $\beta_2$. Let $\mathbf1_v: V(C_n) \to \{0, 1\}$ denote the standard indicator function taking the value $1$ at vertex $v$ and $0$ elsewhere. Now let $\beta_2:= \beta_1 + \gamma \cdot \mathbf1_{u} - \gamma \cdot \mathbf1_{w}$. 

Note that the number of valid initial values for $\beta_1$ is $|Z| = |A| - M(\gamma)$. If $M(\gamma) = |A|$, then $Z = \emptyset$. This implies that no valid $\beta_1$-flow exists at all, meaning $C_n$ violates the fundamental $A$-connectedness property (Definition~\ref{$A$-FRC-def}\ref{FRC-1}), completing the proof immediately. Therefore, we may assume $M(\gamma) < |A|$, which guarantees that $Z$ is non-empty. Consequently, there exists at least one valid $\beta_1$-flow $f_1$, generated by some initial value $z \in Z$.

To convert this valid $f_1$ into a $\beta_2$-flow, we must modify the flow along the path between $u$ and $w$. Suppose we attempt to achieve this by adding $\gamma$ to the edges of $P_1$. The new flow $f_2$ on any edge $e \in E(P_1)$ would algebraically become:
\[ f_2(e) = f_1(e) + \gamma = (z + c_e) + \gamma = (z + \gamma) + c_e. \]
Recall our structural setup: because $z \in Z$, we guaranteed that the shifted value $(z + \gamma)$ falls strictly into the set $X$. Furthermore, because we injectively assigned the value $-x$ to the offsets on $P_1$ for \textit{every} element $x \in X$, there must exist a specific edge $e^* \in E(P_1)$ that received the exact offset $c_{e^*} = -(z+\gamma)$. Evaluating the new flow $f_2$ on this specific edge yields exactly zero:
\[ f_2(e^*) = (z + \gamma) + c_{e^*} = (z + \gamma) - (z + \gamma) = 0_A. \]
This forces a zero value on $e^*$, meaning the conversion along $P_1$ universally fails for all valid initial values $z \in Z$.

Alternatively, modifying $P_2$ requires adding $-\gamma$ to the edges of $P_2$. The new flow on any edge $e \in E(P_2)$ becomes $(z - \gamma) + c_e$. Our construction inherently forces $(z - \gamma) \in Y$, and since the offsets on $P_2$ perfectly cover the additive inverses of $Y$, there must exist an edge $e^{**} \in E(P_2)$ where the flow identically hits zero. 

Since modifying the flow along either path forces an edge to take a value of $0_A$, the flow cannot be converted. Consequently, $C_n$ is not $A$-FRC.
\end{proof}

\end{document}
